\documentclass{article}
\usepackage{amsmath, amssymb}
\usepackage{cancel}
\usepackage{tikz}
\usepackage{mdframed}
\usepackage{hyperref}
\usetikzlibrary{arrows.meta, decorations.markings}

\begin{document}

\section{Lecture 13: Plan for Class}

I missed this day, so these notes will be sparse. The main topic of this lecture was the averaging method. Below are some notes that I had Claude generate from a friend's lecture notes.

\section{Amplitude-Phase Transformation (Variation of Constants)}

\subsection{Motivation}

Many interesting physical systems are \textit{nearly linear} --- they look like a simple harmonic
oscillator with a small perturbation:
\[
    \ddot{x} + \omega^2 x = \varepsilon f(x, \dot{x}, t),
\]
where $\varepsilon \ll 1$ is a small parameter. Examples include:
\begin{itemize}
    \item A pendulum with small damping or driving,
    \item The Mathieu equation (parametric resonance),
    \item Weakly nonlinear springs (Duffing oscillator).
\end{itemize}

When $\varepsilon = 0$, the solution is a pure oscillation:
\[
    x(t) = r_0 \cos(\omega t + \varphi_0),
\]
with \textbf{constant} amplitude $r_0$ and phase $\varphi_0$. The key question is: \textit{what
happens to $r$ and $\varphi$ when $\varepsilon \neq 0$?} Instead of being constants, they will
slowly drift in time. The amplitude-phase method makes this precise.

\subsection{The Ansatz}

We guess that the solution still looks like a cosine, but now with \textbf{slowly varying amplitude
$r(t)$ and phase $\varphi(t)$}:
\[
    x(t) = r(t)\cos(\omega t + \varphi(t)).
\]
This is ``variation of constants'' --- the same idea as in linear ODEs, where formerly-constant
parameters are allowed to vary.

Since we have one function $x(t)$ but two unknowns $r(t)$ and $\varphi(t)$, we have freedom to
impose an extra condition. We choose it to keep the algebra clean by declaring:
\begin{equation}\label{eq:constraint_def}
    \dot{x}(t) = -r(t)\,\omega\sin(\omega t + \varphi(t)). \tag{$*$}
\end{equation}
This says: $\dot{x}$ looks exactly as it would if $r$ and $\varphi$ were constant.

\subsection{The Constraint Equation}

Differentiating $x(t) = r(t)\cos(\omega t + \varphi(t))$ directly via the product and chain rules:
\[
    \dot{x} = \dot{r}\cos(\omega t + \varphi) - r(\omega + \dot{\varphi})\sin(\omega t + \varphi).
\]
But \eqref{eq:constraint_def} demands $\dot{x} = -r\omega\sin(\omega t + \varphi)$. Matching these
two expressions forces:
\[
    \boxed{\dot{r}\cos(\omega t + \varphi) - r\dot{\varphi}\sin(\omega t + \varphi) = 0.}
    \tag{Constraint}
\]
This is the first of two equations needed to determine $\dot{r}$ and $\dot{\varphi}$.

\subsection{The Second Equation from the ODE}

Differentiating \eqref{eq:constraint_def} gives $\ddot{x}$:
\[
    \ddot{x} = -\dot{r}\,\omega\sin(\omega t + \varphi)
               - r\omega(\omega + \dot{\varphi})\cos(\omega t + \varphi).
\]
Substituting into $\ddot{x} + \omega^2 x = \varepsilon f$, the $\omega^2 r\cos(\cdots)$ terms
cancel, leaving:
\[
    \boxed{-\dot{r}\,\omega\sin(\omega t + \varphi) - r\dot{\varphi}\,\omega\cos(\omega t +
    \varphi) = \varepsilon f.} \tag{ODE eq.}
\]

\subsection{Solving for $\dot{r}$ and $\dot{\varphi}$}

We now have a $2\times 2$ linear system in $\dot{r}$ and $\dot{\varphi}$. Solving by multiplying
the constraint by $\omega\sin(\omega t+\varphi)$ and the ODE equation by $\cos(\omega t+\varphi)$,
then adding (and similarly for the other equation), yields the \textbf{exact} averaging form:
\[
    \dot{r} = -\frac{\varepsilon}{\omega}\sin(\omega t + \varphi)\,f(x, \dot{x}, t),
\]
\[
    \dot{\varphi} = -\frac{\varepsilon}{\omega r}\cos(\omega t + \varphi)\,f(x, \dot{x}, t).
\]
These are exact --- no approximation yet. We have simply rewritten the original scalar ODE as a
system for $(r, \varphi)$.

\subsection{Example: The Van der Pol Oscillator}

Consider
\[
    \ddot{x} + x = \varepsilon(1 - x^2)\dot{x},
\]
so $\omega = 1$ and $f(x,\dot{x},t) = (1 - x^2)\dot{x}$. Substituting the ansatz
$x = r\cos\theta$, $\dot{x} = -r\sin\theta$ (where $\theta = t + \varphi$):
\[
    f = \bigl(1 - r^2\cos^2\theta\bigr)(-r\sin\theta).
\]
The exact equations become:
\[
    \dot{r} = \varepsilon r\sin^2\theta\,(1 - r^2\cos^2\theta),
\]
\[
    \dot{\varphi} = \varepsilon\sin\theta\cos\theta\,(1 - r^2\cos^2\theta).
\]
Since $\varepsilon$ is small, $r$ and $\varphi$ change slowly. Over one period, the fast
oscillations average out. Applying the \textbf{method of averaging} --- replacing the right-hand
sides by their averages over one period $T = 2\pi$ while holding $r$ and $\varphi$ fixed:
\[
    \langle \dot{r} \rangle = \frac{\varepsilon r}{2}\!\left(1 - \frac{r^2}{4}\right),
    \qquad
    \langle \dot{\varphi} \rangle = 0.
\]
The $\varphi$ equation says the \textbf{frequency does not shift} to leading order. The $r$
equation has a nontrivial structure:
\begin{itemize}
    \item $r < 2$: $\dot{r} > 0$, amplitude \textbf{grows},
    \item $r > 2$: $\dot{r} < 0$, amplitude \textbf{shrinks},
    \item $r = 2$: $\dot{r} = 0$, amplitude is \textbf{constant}.
\end{itemize}
This predicts a \textbf{stable limit cycle at $r = 2$}, exactly what the Van der Pol oscillator is
known for --- obtained without solving the full nonlinear system.

\subsection{Summary}

\[
    \underbrace{\ddot{x}+\omega^2 x = \varepsilon f}_{\text{nearly linear ODE}}
    \;\xrightarrow{\text{ansatz}}\;
    \underbrace{\dot{r},\,\dot{\varphi}\text{ system}}_{\text{exact}}
    \;\xrightarrow{\text{average}}\;
    \underbrace{\langle\dot{r}\rangle,\,\langle\dot{\varphi}\rangle}_{\text{slow-flow equations}}
\]

The averaged equations for $r$ and $\varphi$ are often much simpler than the original ODE, yet
capture essential long-time behavior: growth or decay of oscillations, limit cycles, and frequency
shifts.

\end{document}





\end{document}